% Generated draft: requires mathematical and source review.
\documentclass[11pt]{article}
\usepackage[margin=25mm]{geometry}
\usepackage{fontspec}
\setmainfont{Latin Modern Roman}
\usepackage{amsmath,amssymb,mathtools,unicode-math}
\setmathfont{Latin Modern Math}
\usepackage{enumitem,needspace,xurl,hyperref}
\setlength{\parindent}{0pt}
\setlength{\parskip}{5pt}
\setlength{\emergencystretch}{3em}
\newcommand{\N}{\mathbb{N}}
\newcommand{\Z}{\mathbb{Z}}
\newcommand{\Q}{\mathbb{Q}}
\newcommand{\R}{\mathbb{R}}
\newcommand{\C}{\mathbb{C}}
\newcommand{\F}{\mathbb{F}}
\newcommand{\A}{\mathbb{A}}
\newcommand{\PP}{\mathbb{P}}
\newcommand{\E}{\mathbb{E}}
\newcommand{\eps}{\varepsilon}
\newcommand{\dd}{\,\mathrm{d}}
\begin{document}
\section*{UnsolvedMath ID: 30005233}
\textbf{Source problem code:} catalogue entry
\subsection{Definitions and mathematical statement}
The supplied record asks for a continuous-time extension of an unspecified discrete-time result for attractive models. Its precise benchmark is the following contact-process statement, recalled as an existing result rather than posed as a new conjecture.

For a lattice dimension \(d\geq 1\), let \(\eta(t))_{t\geq 0}\) be the contact process on \(\{0,1\}^{\Z^d}\). For finite \(\lambda\geq 0\), a site in state \(1\) changes to state \(0\) at rate \(1\), and a site \(x\) in state \(0\) changes to state \(1\) at rate
\[
\lambda\sum_{y:\,\|y-x\|_1=1}\eta_y(t).
\]
Write \(\mathbf 1\) for the configuration identically equal to \(1\), let \(\mathbb P^{\mathbf 1}_\lambda\) denote the law with initial configuration \(\mathbf 1\), and put
\[
\rho_\lambda(t)=\mathbb P^{\mathbf 1}_\lambda\bigl(\eta_0(t)=1\bigr).
\]
The source attributes to Bezuidenhout and Grimmett (1991), as its Theorem 1, the existence of \(\lambda_c\in[0,\infty]\) such that, for finite parameter values,
\[
\begin{aligned}
\lambda>\lambda_c
&\quad\Longrightarrow\quad
\lim_{t\to\infty}\rho_\lambda(t)>0,\\
\lambda<\lambda_c
&\quad\Longrightarrow\quad
\exists\,C,c>0\ \forall\,t\geq0:\quad
\rho_\lambda(t)\leq C e^{-ct}.
\end{aligned}
\]
Here \(C\) and \(c\) may depend on the parameter and the fixed model, but not on \(t\). No assertion at \(\lambda=\lambda_c\) is included. The original excerpt allows \(\lambda\in[0,\infty]\), but does not define the dynamics at \(\lambda=\infty\); the rate construction above therefore specifies only finite \(\lambda\).

The requested extension is to identify a suitable class of finite-range, continuous-time attractive spin systems, containing this contact-process example, and establish an analogous critical-parameter dichotomy. For binary spin systems, attractiveness means preservation of the coordinatewise stochastic order by the evolution; finite range means that a site's transition rates depend only on spins within a fixed finite neighborhood. These properties alone do not specify the intended class or its parameterization. The supplied excerpt does not state the discrete-time theorem to be extended, the general transition rates, the parameter order, the absorbing-state assumptions, or the observable whose exponential decay is requested. Consequently, the displayed contact-process dichotomy is a precise supplied benchmark, while a universally quantified continuous-time extension cannot be formulated precisely from this record alone. In particular, the benchmark concerns occupation at the origin from the all-\(1\) initial configuration; interpreting it as survival from a finite initial configuration for every proposed spin system would require additional hypotheses or an explicitly supplied equivalence.
\subsection{Short English statement}
Extend the supplied discrete-time result to a suitable class of finite-range continuous-time attractive spin systems, recovering the contact-process dichotomy: positive limiting occupation above a critical parameter and exponential decay below it. The general model class and the discrete-time theorem are missing, so the intended extension remains incompletely specified. The contact-process statement is supplied as a known benchmark, and no behavior at criticality is requested.
\subsection{Sources}
\begin{itemize}
\item Supplied local UnsolvedMath registry record, permanent ID \(30005233\), exact source problem code \texttt{catalogue entry}, title ``Continuous-Time Extensions of Attractive Percolation Models,'' in the Oberwolfach Reports Open Problems collection. The record gives proposed year 2022 and names no proposer.
\item Supplied report citation: \emph{Large Scale Stochastic Dynamics (2023)}, Oberwolfach Reports, DOI 10.4171/owr/2022/41, pp. 2399--2466. \url{https://doi.org/10.4171/owr/2022/41}. The excerpt attributes its contact-process Theorem 1 to Bezuidenhout and Grimmett (1991), reference [2], and cites [8, 9] for background; the corresponding bibliography entries are not supplied.
\item Supplied literature source: ``Hartarsky, bootstrap percolation and attractive PCA.'' \url{https://arxiv.org/abs/1806.11405}.
\item Supplied literature source: ``Oberwolfach Report 39/2022.'' \url{https://ems.press/content/serial-article-files/46978}.
\item Literature assessment dated exactly \texttt{2026-08-22}: the record labels the problem open and reports that partial bootstrap-percolation and attractive-PCA results are known, while the requested continuous-time extension and full exponential tail bound remain open in the stated generality. This is supplied dated curation derived from the OWR material, not an independently verified literature assessment; the record itself requests independent verification of time-sensitive status and cited claims.
\item The registry labels the statement \texttt{corrected\_verified} and reports that the official report was checked, the question isolated from surrounding exposition, and notation repaired. These are supplied provenance claims, not verification performed for this draft.
\end{itemize}
\end{document}
